\SGAchapter{XX}{Griffiths 定理}{Le théorème de Griffiths}
\begin{center}N. Katz\end{center}
\SourcePage{351}
\SGAsection{0}{引言}
我们证明 Griffiths \Term{theorem}{定理} [1]。这个\Term{theorem}{定理}给出一种系统的方法（虽然并非构造性的方法），用以“找到”这样的\Term{algebraic-cycle}{代数圈}：它们的 $\mathbf Q_\ell$-\Term{cohomology-class}{上同调类}为零，但其任何倍数都不\Term{algebraic-equivalence}{代数等价}于零。这里给出的证明归功于 Grothendieck，它把 Griffiths 原来多少带有\Term{transcendental}{超越}性质的证明翻译成纯\Term{algebra}{代数}语言。两种证明都以第 XVIII 讲 6.7 的\Term{irreducibility-theorem}{不可约性定理}为关键环节（参见 3.1.6 和 3.2）。

\SGAsection{1}{本原类的形式体系}
以下始终固定一个素数 $\ell$，并假定所考虑的一切\Term{scheme}{概形}的\Term{characteristic}{特征}均不等于 $\ell$。考虑如下情形：
\begin{equation}\tag{1.0.1}
\begin{tikzcd}
X \arrow[rr,"\pi"] \arrow[dr,"f"'] && Y \arrow[dl,"g"]\\
& S &
\end{tikzcd}
\end{equation}
其中作如下假设。
\SGAsourcegap{原书在此直接转到下一页的（1.0.3），未见（1.0.2）；后文 1.2 明确引用（1.0.2），其内容未能从现有原页取得。}
\SourcePage{352}
\noindent（1.0.3）$f,g$ \Term{smooth}{光滑}，\Term{relative-dimension}{相对维数}分别为 $d_1,d_2$；
\noindent（1.0.4）$\pi$ \Term{proper}{固有}。
由于 $X$ 和 $Y$ 在 $S$ 上\Term{smooth}{光滑}，$\pi$ 的\Term{proper}{固有性}给出一个\Term{gysin-morphism}{Gysin 态射}（SGA 5，第 IV 讲\SGAsourcegap{此处原书在“SGA 5 IV”后留空，未给具体引用号。}）：
\begin{equation}\tag{1.0.5}
\pi_*:H^q(X,\mathbf Q_\ell(i))\longrightarrow H^{q-2(d_1-d_2)}(Y,\mathbf Q_\ell(i-(d_1-d_2))).
\end{equation}
在 $S$ 上作\Term{etale}{平展}\Term{localization}{局部化}，还得到 $S$ 上的\Term{sheaf}{层}\Term{morphism}{态射}
\begin{equation}\tag{1.0.6}
\pi_*:R^qf_*\mathbf Q_\ell(i)\longrightarrow R^{q-2(d_1-d_2)}g_*\mathbf Q_\ell(i-d_1-d_2),
\end{equation}
再施以\Term{functor}{函子} $H^p(S,-)$，得到\Term{morphism}{态射}
\begin{equation}\tag{1.0.7}
\pi_*:H^p(S,R^qf_*\mathbf Q_\ell(i))\longrightarrow H^p(S,R^{q-2(d_1-d_2)}g_*\mathbf Q_\ell(i-d_1+d_2)).
\end{equation}
\begin{SGAresult}{定义}{1.0.8}
设 $(E_r^{p,q},d_r^{p,q})$ 和 $('E_r^{p,q},'d_r^{p,q})$ 是两个\Term{spectral-sequence}{谱序列}（取值于一个给定的\Term{abelian-category}{阿贝尔范畴}），且 $(a,b)\in\mathbf Z\times\mathbf Z$。$('E_r^{p,q},'d_r^{p,q})$ 按 $(a,b)$ 的\Term{shift}{平移}，是指\Term{spectral-sequence}{谱序列} $(''E_r^{p,q},''d_r^{p,q})$，其中
\begin{equation}\tag{1.0.9}
\begin{cases}
''E_r^{p,q}='E_r^{p+a,q+b},\\
''d_r^{p,q}=d_r^{p+a,q+b}.
\end{cases}
\end{equation}
\end{SGAresult}
\SourcePage{353}
\textbf{1.0.10。} 从 $(E_r^{p,q},d_r^{p,q})$ 到 $(''E_r^{p,q},''d_r^{p,q})$ 的\Term{spectral-sequence}{谱序列}\Term{morphism}{态射}，称为从 $(E_r^{p,q},d_r^{p,q})$ 到 $('E_r^{p,q},'d_r^{p,q})$ 的\Term{bidegree}{双次数}为 $(a,b)$ 的\Term{morphism}{态射}。

以下\Term{lemma}{引理}是\Term{gysin-morphism}{Gysin 态射}理论（SGA 5，第 IV 讲）的直接推论。
\begin{SGAresult}{\Term{lemma}{引理}}{1.1}
存在以下两个 Leray \Term{spectral-sequence}{谱序列}之间的\Term{canonical}{典范}\Term{morphism}{态射} $\pi_*$，其\Term{bidegree}{双次数}为 $(a,-2(d_1-d_2))$：
\begin{equation}\tag{1.1.0}
R^p\alpha_*R^qf_*\mathbf Q_\ell(i)\Longrightarrow R^{p+q}(\alpha f)_*\mathbf Q_\ell(i)
\end{equation}
和
\begin{equation}\tag{1.1.1}
R^p\alpha_*R^qg_*\mathbf Q_\ell(i-d_1+d_2)\Longrightarrow R^{p+q}(\alpha g)_*\mathbf Q_\ell(i-d_1+d_2).
\end{equation}
它在 $E_2$ 项上由（1.0.7）给出，在\Term{abutment}{终项}上由（1.0.5）给出。
\end{SGAresult}
\textbf{1.2。} 现在考虑情形
\begin{equation}\tag{1.2.1}
\begin{tikzcd}
& X\times_S Y \arrow[dl,"\operatorname{pr}_1"'] \arrow[dr,"\operatorname{pr}_2"] &\\
X \arrow[dr,"f"'] && Y \arrow[dl,"g"]\\
& S \arrow[d,"\alpha"] &\\
& T &
\end{tikzcd}
\end{equation}
\SourcePage{354}
其中 $S,X,f,Y,g$ 如（1.0.2）—（1.0.4），但现在还假定 $f$ \Term{proper}{固有}。设
\begin{equation}\tag{1.2.2}
z\in H^a(X\times_S Y,\mathbf Q_\ell(b)).
\end{equation}
则 $z$ 诱导一个仍记为 $z$ 的映射
\begin{equation}\tag{1.2.3}
z:H^q(X,\mathbf Q_\ell(i))\longrightarrow H^{q+a-2d_1}(Y,\mathbf Q_\ell(i+b-d_1)),
\end{equation}
其定义公式为
\begin{equation}\tag{1.2.4}
z(x)=\operatorname{pr}_{2*}(z\wedge\operatorname{pr}_1^*(x)),
\end{equation}
其中 $\operatorname{pr}_{2*}$ 是\Term{gysin-morphism}{Gysin 同态}。在 $S$ 上作\Term{etale}{平展}\Term{localization}{局部化}，$z$ 还定义了
\begin{equation}\tag{1.2.5}
z:R^qf_*\mathbf Q_\ell(i)\longrightarrow R^{q+a-2d_1}g_*\mathbf Q_\ell(i+b-d_1),
\end{equation}
再施以 $H^p(S,-)$，得到映射
\begin{equation}\tag{1.2.6}
H^p(S,R^qf_*\mathbf Q_\ell(i))\longrightarrow H^p(S,R^{q+a-2d_1}g_*\mathbf Q_\ell(i+b-d_1)).
\end{equation}
\begin{SGAresult}{\Term{lemma}{引理}}{1.3}
类 $z$ 定义以下两个 Leray \Term{spectral-sequence}{谱序列}之间\Term{bidegree}{双次数}为 $(0,b-2d_1)$ 的\Term{morphism}{态射}：
\begin{equation}\tag{1.2.7}
H^p(S,R^qf_*\mathbf Q_\ell(i))\Longrightarrow H^{p+q}(X,\mathbf Q_\ell(i))
\end{equation}
和
\begin{equation}\tag{1.2.8}
H^p(S,R^qg_*\mathbf Q_\ell(i+b-d_1))\Longrightarrow H^{p+q}(Y,\mathbf Q_\ell(i+b-d_1)).
\end{equation}
\SourcePage{355}
该\Term{morphism}{态射}在 $E_2$ 项上由（1.2.6）给出，在\Term{abutment}{终项}上由（1.2.3）给出。
\end{SGAresult}
\begin{proof}
由公式（1.2.4），并注意到“乘以 $z$”定义了以下两个 Leray \Term{spectral-sequence}{谱序列}之间\Term{bidegree}{双次数}为 $(0,a)$ 的\Term{morphism}{态射}：
\begin{equation}\tag{1.2.9}
H^p(S,R^q(\operatorname{pr}_1 f)_*\mathbf Q_\ell(i))\Longrightarrow H^{p+q}(X\times_S Y,\mathbf Q_\ell(i)),
\end{equation}
\begin{equation}\tag{1.2.10}
H^p(S,R^{q+a}(\operatorname{pr}_1 f)_*\mathbf Q_\ell(i+b))\Longrightarrow H^{p+q+a}(X\times_S Y,\mathbf Q_\ell(i+b)),
\end{equation}
只需将 1.1 应用于情形
\begin{equation}\tag{1.2.11}
\begin{tikzcd}
X\times_S Y \arrow[rr,"\operatorname{pr}_2"] \arrow[dr,"f\operatorname{pr}_1"'] && Y \arrow[dl,"g"]\\
& S &
\end{tikzcd}
\end{equation}
即可。
\end{proof}
\begin{SGAresult}{推论}{1.4}
在（1.2.1）—（1.2.2）的情形下，\Term{morphism}{态射}（1.2.3）
\begin{equation}\tag{1.4.1}
z:H^q(X,\mathbf Q_\ell(i))\longrightarrow H^{q+a-2d_1}(Y,\mathbf Q_\ell(i+b-d_1))
\end{equation}
把\Term{primitive-class}{本原类}（第 XVIII 讲 5.8.1）映为\Term{primitive-class}{本原类}，且以下图表交换：
\SourcePage{356}
\begin{equation}\tag{1.4.2}
\begin{tikzcd}[column sep=small]
\operatorname{Prim}^q(X/S,\mathbf Q_\ell(i)) \arrow[r,"z"] \arrow[d,"u_f"'] & \operatorname{Prim}^{q+a-2d_1}(Y/S,\mathbf Q_\ell(i+b-d_1)) \arrow[d,"u_g"]\\
H^1(S,R^{q-1}f_*\mathbf Q_\ell(i)) \arrow[r,"z"] & H^1(S,R^{q+a-2d_1-1}g_*\mathbf Q_\ell(i+b-d_1)).
\end{tikzcd}
\end{equation}
其中 $u_f,u_g$ 是第 XVIII 讲 5.8.1 中的\Term{canonical}{典范}\Term{homomorphism}{同态}。
\end{SGAresult}
\textbf{1.5。} 设 $\bar\eta$ 是 $S$ 的一个\Term{geometric-point}{几何点}（位于 $S$ 的某一点 $\eta$ 之上），并设
\begin{equation}\tag{1.5.1}
z_{\bar\eta}\in H^a(X_{\bar\eta}\times_{k(\bar\eta)}Y_{\bar\eta},\mathbf Q_\ell(b))
\end{equation}
为 $z$ 在
\begin{equation}\tag{1.5.2}
X\times_S Y\longrightarrow S
\end{equation}
的相应\Term{geometric-fibre}{几何纤维}上诱导的类。
\begin{SGAresult}{命题}{1.5.3}
假定 $f,g$ \Term{proper}{固有}，$S$ \Term{connected}{连通}，并设 $q\in\mathbf Z$。若类 $z_{\bar\eta}$ 诱导零\Term{homomorphism}{同态}
\begin{equation}\tag{1.5.3}
0=z_{\bar\eta}:H^{q-1}(X_{\bar\eta},\mathbf Q_\ell(i))\longrightarrow H^{q+a-2d_1-1}(Y_{\bar\eta},\mathbf Q_\ell(i+b-d_1)),
\end{equation}
则对所有 $x\in\operatorname{Prim}^q(X/S,\mathbf Q_\ell(i))$，都有
\begin{equation}\tag{1.5.4}
u_g(z(x))=0\quad\text{于 }H^1(S,R^{q+a-2d_1-1}g_*\mathbf Q_\ell(i+b-d_1)).
\end{equation}
\end{SGAresult}
\SourcePage{357}
\begin{proof}
由（1.4.2）的\Term{commutativity}{交换性}，只需证明 $z$ 诱导零\Term{homomorphism}{同态}
\begin{equation}\tag{1.5.5}
z:R^{q-1}f_*\mathbf Q_\ell(i)\xrightarrow{\ 0\ }R^{q+a-2d_1-1}g_*\mathbf Q_\ell(i+b-d_1).
\end{equation}
而 $f,g$ 均\Term{proper}{固有}且\Term{smooth}{光滑}，所以\Term{sheaf}{层} $R^\bullet f_*\mathbf Q_\ell(\bullet)$ 和 $R^\bullet g_*\mathbf Q_\ell(\bullet)$ 都是\Term{twisted-constant}{扭常值}的（SGA 4，第 XVI 讲 2.2），因而\Term{equivalence}{等价}于它们在 $\bar\eta$ 处的值，并把这些值看作 $\pi_1(S,\bar\eta)$-模。
\end{proof}
\textbf{1.6。Manin 的一个构造。}
假定给定 $f:X\to S$ 的两个\Term{section}{截面}：
\begin{equation}\tag{1.6.0}
\begin{tikzcd}
S \arrow[r,bend left=35,"t_1"] \arrow[r,bend right=35,"t_2"'] & X \arrow[l,"f" description]
\end{tikzcd}
\end{equation}
对每个\Term{section}{截面}，赋予其\Term{cohomology-class}{上同调类}
\begin{equation}\tag{1.6.1}
\operatorname{cl}(t_i)=t_{i*}(1)\in H^{2d_1}(X,\mathbf Q_\ell(d_1)).
\end{equation}
若 $f:X\to S$ 的\Term{fibre}{纤维}\Term{geometrically-connected}{几何连通}（更\Term{generic}{一般}地，若对 $s\in S$，$f(s)$ 和 $g(s)$ 属于 $X_s$ 的同一个\Term{connected-component}{连通分支}），则显然
\begin{equation}\tag{1.6.2}
\operatorname{cl}(t_1)-\operatorname{cl}(t_2)\in\operatorname{Prim}^{2d_1}(X/S,\mathbf Q_\ell(d_1)).
\end{equation}
\SourcePage{358}
用
\begin{equation}\tag{1.5.3}
u_f(t_1-t_2)\in H^1(S,R^{2d_1-1}f_*\mathbf Q_\ell(d_1))
\end{equation}
表示它在
\begin{equation}\tag{1.6.4}
u_f:\operatorname{Prim}^{2d_1}(X/S,\mathbf Q_\ell(d_1))\longrightarrow H^1(S,R^{2d_1-1}f_*\mathbf Q_\ell(d_1))
\end{equation}
下的像。
（Manin 在 [5] 中研究的是这个元素的一个类似物 $\mu(t_1-t_2)$。）

\textbf{1.6.5。} 现在设 $z$ 是 $X\times_S Y$ 上的一个\Term{algebraic-cycle}{代数圈}，在 $X\times_S Y$ 的每条\Term{fibre}{纤维}上的\Term{codimension}{余维}均 $\geq b$\footnote{若 $S$ 在某个\Term{scheme}{概形} $T$ 上\Term{smooth}{光滑}（例如 $T=\operatorname{Spec}k$，其中 $k$ 是域），则只需假定 $z$ 在 $T$ 上的每条\Term{fibre}{纤维}上的\Term{codimension}{余维}均 $\geq b$。}，它确定一个\Term{cohomology-class}{上同调类}（SGA 5，第 IV 讲\SGAsourcegap{原书此处的“SGA 5 IV”之后未填写具体引用号。}）：
\begin{equation}\tag{1.6.6}
\operatorname{cl}(z)\in H^{2b}(X\times_S Y,\mathbf Q_\ell(b)).
\end{equation}
这个类定义一个映射（参见（1.2.3））
\begin{equation}\tag{1.6.7}
\operatorname{cl}(z):H^{2d_1}(X,\mathbf Q_\ell(d_1))\longrightarrow H^{2b}(Y,\mathbf Q_\ell(b)),
\end{equation}
而且，若圈 $z(t_1)$ 和 $z(t_2)$ 有定义，则由\Term{cup-product}{杯积}与\Term{cycle-intersection}{圈的交}的相容性（SGA 5，第 IV 讲\SGAsourcegap{原书关于\Term{cup-product}{杯积}与\Term{cycle-intersection}{圈的交}相容性的“SGA 5 IV”引注之后留空。}），有
\SourcePage{359}
\begin{equation}\tag{1.6.8}
\operatorname{cl}(z)(\operatorname{cl}(t_1)-\operatorname{cl}(t_2))=\operatorname{cl}(z(t_1)-z(t_2)).
\end{equation}
（1.4.2）的\Term{commutativity}{交换性}给出
\begin{equation}\tag{1.6.9}
\operatorname{cl}(z)(u_f(t_1-t_2))=u_g(\operatorname{cl}(z(t_1)-z(t_2))).
\end{equation}
于是，由 1.5.3 得到以下推论。
\begin{SGAresult}{命题}{1.7}
设 $t_1,t_2$ 是 $f:X\to S$ 的两个\Term{section}{截面}，假定 $f$ 的\Term{fibre}{纤维}\Term{connected}{连通}；并设 $z$ 是 $X\times_S Y$ 上\Term{codimension}{余维}为 $b$ 的\Term{algebraic-cycle}{代数圈}，使\Term{algebraic-cycle}{代数圈}
\begin{equation}\tag{1.7.1}
z(t_1)-z(t_2)
\end{equation}
有定义。则它的\Term{cohomology-class}{上同调类}是\Term{primitive-class}{本原的}（第 XVIII 讲 5.8.1）。若 $X_{\bar\eta}\times Y_{\bar\eta}$ 上的\Term{algebraic-cycle}{代数圈} $z_{\bar\eta}$ 诱导零\Term{homomorphism}{同态}
\begin{equation}\tag{1.7.2}
z_{\bar\eta}:H^{2d_1-1}(X_{\bar\eta},\mathbf Q_\ell(d_1))\xrightarrow{\ 0\ }H^{2b-1}(Y_{\bar\eta},\mathbf Q_\ell(b)),
\end{equation}
则
\begin{equation}\tag{1.7.3}
u_g(z(t_1)-z(t_2))=0\quad\text{于 }H^1(S,R^{2b-1}g_*\mathbf Q_\ell(b)),
\end{equation}
其中 $u_g$ 是第 XVIII 讲 5.8.1 中出现于 1.4 的\Term{homomorphism}{同态}。
\end{SGAresult}

\SourcePage{360}
\SGAsection{2}{关于层次的回顾（参见 [2]，9.7、10.1）}
\textbf{2.0.1。} 设 $X$ 在\Term{algebraically-closed}{代数闭}域 $k$ 上\Term{proper}{固有}且\Term{smooth}{光滑}。由（1.2.3）可知，对任意
\noindent（2.0.2）$T$ 在 $k$ 上\Term{proper}{固有}、\Term{smooth}{光滑}且\Term{connected}{连通}，\Term{dimension}{维数}为 $d$；
以及任意
\noindent（2.0.3）$T\times X$ 上\Term{codimension}{余维}为 $b$ 的\Term{algebraic-cycle}{代数圈} $z$；
\Term{cohomology-class}{上同调类}
\begin{equation}\tag{2.0.4}
\operatorname{cl}(z)\in H^{2b}(T\times X,\mathbf Q_\ell(b))
\end{equation}
给出映射
\begin{equation}\tag{2.0.5}
\operatorname{cl}(z):H^q(T,\mathbf Q_\ell(i))\longrightarrow H^{q+2b-2d}(X,\mathbf Q_\ell(i+b-d)).
\end{equation}
定义 $H^q(X,\mathbf Q_\ell(i))$ 的\Term{coniveau-filtration}{余层次滤过}（filtration par le co-niveau）$(N^j)$ 如下：
\begin{equation}\tag{2.0.6}
N^jH^q(X,\mathbf Q_\ell(i))=\sum_{(T,z)}\operatorname{Im}\bigl(\operatorname{cl}(z):H^{q-2j}(T,\mathbf Q_\ell(i-j))\longrightarrow H^q(X,\mathbf Q_\ell(i))\bigr),
\end{equation}
其中求和遍历所有满足下列条件的二元组 $(T,z)$：
\noindent（2.0.7）$T/k$ \Term{proper}{固有}、\Term{smooth}{光滑}且\Term{connected}{连通}，\Term{dimension}{维数}为 $d$；$z$ 是 $T\times X$ 上\Term{codimension}{余维}为 $j+d$ 的\Term{algebraic-cycle}{代数圈}。
应当指出，有时也使用递增的\Term{level-filtration}{层次滤过}（filtration par le niveau）$F^j$，定义为
\begin{equation}\tag{2.0.8}
F^jH^q(X,\mathbf Q_\ell(i))=N^{[\frac{q-j}{2}]}H^q(X,\mathbf Q_\ell(i)).
\end{equation}
\SourcePage{361}
把 $(T,z)$ 换成 $(\mathbf P^1\times T,\{\mathrm{pt}\}\times z)$，便可看出
\begin{equation}\tag{2.0.9}
N^jH^q(X,\mathbf Q_\ell(i))\subset N^{j-1}H^q(X,\mathbf Q_\ell(i)).
\end{equation}
显然有：
\begin{SGAresult}{命题}{2.1}
设 $X$ 在域 $k$ 上\Term{proper}{固有}、\Term{smooth}{光滑}且\Term{geometrically-connected}{几何连通}；$\bar k$ 是 $k$ 的一个\Term{algebraic-closure}{代数闭包}，$X_{\bar k}=X\otimes_k\bar k$。则 $\operatorname{Gal}(\bar k/k)$ 在 $H^q(X_{\bar k},\mathbf Q_\ell(i))$ 上的\Term{action}{作用}保持\Term{coniveau-filtration}{余层次滤过}。
\end{SGAresult}
标准论证（参见 [2]）给出以下命题。
\begin{SGAresult}{命题}{2.2}
\Term{level-filtration}{层次滤过}在\Term{algebraically-closed}{代数闭}\Term{field-extension}{域扩张} $k'/k$ 下\Term{invariant}{不变}。
\end{SGAresult}
\begin{SGAresult}{命题}{2.3}
设 $S$ 是一个\Term{reduced}{约化}\Term{irreducible}{不可约}\Term{scheme}{概形}，其\Term{generic-point}{泛点}为 $\eta$，并设
\[\pi:X\longrightarrow S\]
是一个\Term{proper}{固有}\Term{smooth}{光滑}\Term{morphism}{态射}，具有\Term{geometrically-connected}{几何连通}\Term{fibre}{纤维}。若给定整数 $j,q$，对\Term{geometric-generic-fibre}{几何泛纤维} $X_{\bar\eta}$ 有
\begin{equation}\tag{2.3.1}
H^q(X_{\bar\eta},\mathbf Q_\ell)=N^jH^q(X_{\bar\eta},\mathbf Q_\ell),
\end{equation}
则存在一个\Term{nonempty}{非空}\Term{open-subscheme}{开子概形} $U\subset S$，使得对每点 $u\in U$，相应的\Term{geometric-fibre}{几何纤维} $X_{\bar u}$ 都满足
\begin{equation}\tag{2.3.2}
H^q(X_{\bar u},\mathbf Q_\ell)=N^jH^q(X_{\bar u},\mathbf Q_\ell).
\end{equation}
\end{SGAresult}

\SourcePage{362}
\SGAsection{3}{对 Lefschetz 铅笔的应用}
\textbf{3.1。} 从现在起，固定：

\noindent（3.1.0）一个\Term{algebraically-closed}{代数闭}域 $k$；

\noindent（3.1.1）一个 $k$ 上\Term{projective}{射影}、\Term{smooth}{光滑}且\Term{connected}{连通}的\Term{scheme}{概形} $X\subset\mathbf P^r$，其\Term{dimension}{维数} $n$ 为偶数，并满足（LV）（第 XVIII 讲 5.2.2）。

取

\noindent（3.1.2）一个由\Term{degree}{次数} $d$ 的\Term{hypersurface}{超曲面}组成的\Term{lefschetz-pencil}{Lefschetz 铅笔} $D=\{H_t\}_{t\in\mathbf P^1}$，其\Term{generic-fibre}{泛纤维}为
\noindent（3.1.3）$X_\eta$，其中 $\eta$ 是 $\mathbf P^1$ 的\Term{generic-point}{泛点}。
置（参见第 XVIII 讲 5.4.4）
\begin{equation}\tag{3.1.4}
m:X_{\bar\eta}\hookrightarrow X\times_k k(\bar\eta)
\end{equation}
为\Term{inclusion}{包含映射}，并令
\begin{equation}\tag{3.1.5}
E^{n-1}(X_{\bar\eta},\mathbf Q_\ell(i))=\bigl(m^*H^{n-1}(X,\mathbf Q_\ell(i))\bigr)^\perp\subset H^{n-1}(X_{\bar\eta},\mathbf Q_\ell(i)).
\end{equation}
我们已经看到（第 XVIII 讲 6.7.2），\Term{picard-lefschetz-formula}{Picard--Lefschetz 公式}蕴含：

\noindent（3.1.6）若 $n$ 为偶数，则 $\operatorname{Gal}(\overline{k(\eta)}/k(\eta))$ 的每个\Term{finite-index-subgroup}{有限指数子群}都在 $E^{2n-1}(X_{\bar\eta},\mathbf Q_\ell(i))$ 上\Term{irreducible-representation}{不可约地作用}。

\begin{SGAresult}{命题}{3.2}
设 $D=\{H_t\}$ 是一个\Term{lefschetz-pencil}{Lefschetz 铅笔}，并假定 $n=2m$ 为偶数。则或者
\SourcePage{363}
\begin{equation}\tag{3.2.1}
H^{n-1}(X_{\bar\eta},\mathbf Q_\ell(i))\subset m^*H^{n-1}(X,\mathbf Q_\ell(i))+N^{m-1}H^{n-1}(X_{\bar\eta},\mathbf Q_\ell(i)),
\end{equation}
或者有：

\noindent（3.2.2）对 $k(\bar\eta)$ 上的每条\Term{proper}{固有}\Term{smooth}{光滑}\Term{curve}{曲线} $C$，以及 $C\times X_{\bar\eta}$ 上每个\Term{codimension}{余维}为 $m$ 的\Term{algebraic-cycle}{代数圈} $z$，诱导的\Term{homomorphism}{同态}（1.2.3）
\begin{equation}\tag{3.2.3}
\operatorname{cl}(z):H^1(C,\mathbf Q_\ell(1+i-m))\longrightarrow H^{2m-1}(X_{\bar\eta},\mathbf Q_\ell(i))
\end{equation}
的像包含于 $m^*H^{2m-1}(X,\mathbf Q_\ell(i))$ 中。
\end{SGAresult}
\begin{proof}
假定某个这样的\Term{homomorphism}{同态}（3.2.3）的像不包含于 $m^*H^{n-1}(X,\mathbf Q_\ell(i))$ 中；\Term{equivalence}{等价}地，它在
\[
E^{n-1}(X_{\bar\eta},\mathbf Q_\ell(i))=H^{n-1}(X_{\bar\eta},\mathbf Q_\ell(i))/H^{n-1}(X,\mathbf Q_\ell(i))
\]
中的像非零。这个像显然在 $\operatorname{Gal}(\overline{k(\eta)}/k(\eta))$ 的某个\Term{finite-index-subgroup}{有限指数子群}的\Term{action}{作用}下\Term{stable}{稳定}，因此 $E^{n-1}(X_{\bar\eta},\mathbf Q_\ell(i))$ 的\Term{irreducibility}{不可约性}（3.1.6）可改写为
\[
H^{n-1}(X_{\bar\eta},\mathbf Q_\ell(i))=m^*H^{n-1}(X,\mathbf Q_\ell(i))+\operatorname{Im}(3.2.3).
\]
而由定义，
\[
\begin{aligned}
\operatorname{Im}(3.2.3)&=\operatorname{cl}(z)\bigl(H^1(C,\mathbf Q_\ell(1+i-m))\bigr)\\
&\subset N^{m-1}H^{n-1}(X_{\bar\eta},\mathbf Q_\ell(i)).
\end{aligned}
\]
所以确实得到
\[
H^{n-1}(X_{\bar\eta},\mathbf Q_\ell(i))=m^*H^{n-1}(X,\mathbf Q_\ell(i))+N^{m-1}H^{n-1}(X_{\bar\eta},\mathbf Q_\ell(i)).
\]
\end{proof}
\SourcePage{364}
\textbf{3.3。条件（B）。}

\textbf{3.3.1。} 称\Term{inclusion}{包含映射}
\[m:X_{\bar\eta}\hookrightarrow X\times_k k(\bar\eta)\]
满足条件（B），如果存在

\noindent（3.3.2）$X_{\bar\eta}\times(X\times_k k(\bar\eta))$ 上\Term{codimension}{余维}为 $n-1$ 的\Term{algebraic-cycle}{代数圈} $V_{\bar\eta}$，使它定义的映射
\begin{equation}\tag{3.3.3}
V_{\bar\eta}:H^{n-1}(X_{\bar\eta},\mathbf Q_\ell(i))\longrightarrow H^{n-1}(X,\mathbf Q_\ell(i))
\end{equation}
是
\begin{equation}\tag{3.3.4}
m^*:H^{n-1}(X,\mathbf Q_\ell(i))\longrightarrow H^{n-1}(X_{\bar\eta},\mathbf Q_\ell(i))
\end{equation}
的\Term{left-inverse}{左逆}，即满足
\begin{equation}\tag{3.3.5}
V_{\bar\eta}m^*=\operatorname{id}.
\end{equation}
\textbf{3.4。} 如果\Term{operator}{算子} $\Lambda$ [4，1.4.2.1] 是\Term{algebra}{代数}的，则条件（B）成立（参见 [4，2.12 (ii)]）。人们猜想 $\Lambda$ 总是\Term{algebra}{代数}的，因此条件（B）总成立；而且已知在许多情形中它成立 [4]。当 $X$ 是偶数维\Term{complete-intersection}{完全交}时，这个条件显然成立，因为这时已知（第 IX 讲 1.6）
\begin{equation}\tag{3.4.1}
H^{n-1}(X,\mathbf Q_\ell(i))=0.
\end{equation}

\SourcePage{365}
\SGAsection{4}{Griffiths 定理}
\begin{SGAresult}{\Term{theorem}{定理}}{4.1}
设 $X$ 在 $k$ 上\Term{projective}{射影}、\Term{smooth}{光滑}且\Term{connected}{连通}，其\Term{dimension}{维数}为偶数 $n=2m\geq4$。设 $D$ 是由\Term{degree}{次数} $d$ 的\Term{hypersurface}{超曲面}\Term{section}{截面}组成的\Term{lefschetz-pencil}{Lefschetz 铅笔}。假定：

\noindent（4.1.0）$X$ 满足（LV）（第 XVIII 讲 5.2.2）；

\noindent（4.1.1）\Term{lefschetz-pencil}{铅笔} $D$ 满足假设（A）（第 XVIII 讲 5.3）；

\noindent（4.1.2）\Term{lefschetz-pencil}{铅笔} $D$ 满足（3.2.2），即（3.2）\Term{vanishing-cohomology}{消失上同调} $E^{n-1}(X_{\bar\eta},\mathbf Q_\ell)$ 的\Term{level}{层次}并非 $\leq1$；

\noindent（4.1.3）\Term{inclusion}{包含映射} $m^*:X_{\bar\eta}\hookrightarrow X\times_k k(\bar\eta)$ 满足假设（B）（3.3）。

若 $z$ 是 $X$ 上\Term{codimension}{余维}为 $m$ 的\Term{algebraic-cycle}{代数圈}，它作为\Term{algebraic-cycle}{代数圈}在 $X_{\bar\eta}$ 上的\Term{restriction}{限制}满足
\noindent（4.1.4）$z_{\bar\eta}=z|X_{\bar\eta}$ \Term{algebraic-equivalence}{代数等价}于零，
则

\noindent（4.1.5）$z$ 的 $\mathbf Q_\ell$-\Term{cohomology-class}{上同调类}为零，即 $\operatorname{cl}(z)\in H^{2m}(X,\mathbf Q_\ell(m))$ 为零。
\end{SGAresult}
\textbf{4.2。关于 $X$ 为\Term{complete-intersection}{完全交}的情形的一点说明。}
我们已指出（参见第 XVIII 讲 5.2.5 及（3.4）），此时假设（4.1.0）和（4.1.3）自动满足。至于假设（4.1.1），我们已看到（第 XVIII 讲 6.3.4、6.4），由\Term{degree}{次数} $d>2n$ 的\Term{hypersurface}{超曲面}组成的\Term{lefschetz-pencil}{铅笔}满足它。另一方面，下一讲将对由\Term{degree}{次数} $d>2n$ 的\Term{hypersurface}{超曲面}组成的“\Term{generic}{一般}”\Term{lefschetz-pencil}{铅笔}证明条件（4.1.2）（第 XXI 讲 5.2）。所以，由\Term{degree}{次数} $d>2n$ 的\Term{hypersurface}{超曲面}组成的“\Term{generic}{一般}”\Term{lefschetz-pencil}{铅笔}满足\Term{theorem}{定理} 4.1 的全部假设。

\SourcePage{366}
\textbf{4.3。} 回到任意 $X$ 的情形，一个很可信的看法是：当\Term{degree}{次数} $d$ \Term{large-relative-to-cardinal}{足够大}时，$D$ 满足 4.1 的全部条件，至少在 $D$ 足够\Term{generic}{一般}时如此。对（4.1.1），我们已经知道确实如此（第 XVIII 讲 6.4）；并且已经指出，人们猜想（4.1.0）和（4.1.3）在任何情形下都成立（当 $\operatorname{char}k=0$ 时，（4.1.0）确实总成立）。至于条件（4.1.2），当 $\operatorname{char}k=0$ 时，由 \Term{hodge-theory}{Hodge 理论}知它在 $d$ \Term{large-relative-to-cardinal}{足够大}时成立；而一些启发性的论证（读者可参见下一讲的附录）使人相信，在 $\operatorname{char}k>0$ 时仍然如此。更\Term{generic}{一般}地，这些论证（\Term{hodge-theory}{Hodge 理论}，以及相应的\Term{crystalline-cohomology}{晶体上同调}）证明了（相应地，使人相信）：对一个 $n$ 维\Term{smooth}{光滑}\Term{projective}{射影}\Term{variety}{簇} $X$，以及由\Term{degree}{次数} $d$ \Term{large-relative-to-cardinal}{足够大}的\Term{hypersurface}{超曲面}组成的任意\Term{lefschetz-pencil}{Lefschetz 铅笔} $D$（相应地，由\Term{degree}{次数} $d$ 的\Term{hypersurface}{超曲面}组成的任意“足够\Term{generic}{一般}”的\Term{lefschetz-pencil}{铅笔} $D$），\Term{vanishing-cohomology}{消失上同调} $E^{n-1}(X_{\bar\eta},\mathbf Q_\ell)$ 的\Term{level}{层次}等于 $n-1$（即其\Term{level}{层次}并非 $\leq n-3$，亦即其\Term{coniveau}{余层次}并非 $1$）。

\textbf{4.4。Griffiths \Term{theorem}{定理}的证明。}
设 $z$ 是 $X$ 上\Term{codimension}{余维}为 $n$ 的\Term{algebraic-cycle}{代数圈}，满足
\noindent（4.4.1）$z_{\bar\eta}=z|X_{\bar\eta}$ \Term{algebraic-equivalence}{代数等价}于零。
由定义，存在：

\noindent（4.4.2）一条 $k(\bar\eta)$ 上\Term{proper}{固有}、\Term{smooth}{光滑}且\Term{connected}{连通}的\Term{curve}{曲线} $C_{\bar\eta}$；
\SourcePage{367}
\noindent（4.4.3）两个\Term{section}{截面} $t_i(\bar\eta):\operatorname{Spec}(k(\bar\eta))\longrightarrow C$，$i=1,2$；
\noindent（4.4.4）$C_{\bar\eta}\times X_{\bar\eta}$ 上一个\Term{codimension}{余维}为 $n$ 的\Term{algebraic-cycle}{代数圈} $P_{\bar\eta}$，

使得
\begin{equation}\tag{4.4.5}
z_{\bar\eta}=P_{\bar\eta}(t_1(\bar\eta))-P_{\bar\eta}(t_2(\bar\eta))\quad\text{于 }X_{\bar\eta}.
\end{equation}
把圈 $P_{\bar\eta}$ 看作 $(C_{\bar\eta},X_{\bar\eta})$ 之间的\Term{algebraic-correspondence}{代数对应}，它定义一个仍记为 $P_{\bar\eta}$ 的\Term{cohomological-correspondence}{上同调对应}，也就是一个\Term{homomorphism}{同态}
\begin{equation}\tag{4.4.6}
P_{\bar\eta}:H^*(C_{\bar\eta},\mathbf Q_\ell(i))\longrightarrow H^{*+2(m-1)}(X_{\bar\eta},\mathbf Q_\ell(i+m-1)).
\end{equation}
类似地，考虑由（4.1.3）的假设（B）给出的 $(X_{\bar\eta},X\otimes_k k(\bar\eta))$ 之间的\Term{algebraic-correspondence}{代数对应} $V_{\bar\eta}$（3.3.2），它定义
\begin{equation}\tag{4.4.7}
V_{\bar\eta}:H^*(X_{\bar\eta},\mathbf Q_\ell(i+m-1))\longrightarrow H^*(X\otimes_k k(\bar\eta),\mathbf Q_\ell(i+m-1)).
\end{equation}
于是得到 $(C_{\bar\eta},X\otimes_k k(\bar\eta))$ 之间的复合\Term{algebraic-correspondence}{代数对应} $Q_{\bar\eta}=V_{\bar\eta}\circ P_{\bar\eta}$，它定义
\begin{equation}\tag{4.4.8}
Q_{\bar\eta}=V_{\bar\eta}\circ P_{\bar\eta}:H^*(C_{\bar\eta},\mathbf Q_\ell(i))\longrightarrow H^{*+2(m-1)}(X\otimes_k k(\bar\eta),\mathbf Q_\ell(i+m-1)).
\end{equation}
再考虑\Term{inclusion}{包含映射}
\[m:X_{\bar\eta}\hookrightarrow X\otimes_k k(\bar\eta),\]
它给出 $(X\otimes_k k(\bar\eta),X_{\bar\eta})$ 之间的\Term{algebraic-correspondence}{代数对应} $m^*$，从而给出 $(C_{\bar\eta},X_{\bar\eta})$ 之间的复合\Term{algebraic-correspondence}{对应} $m^*Q_{\bar\eta}$，其\Term{action}{作用}为
\SourcePage{368}
\begin{equation}\tag{4.4.9}
m^*Q_{\bar\eta}=m^*V_{\bar\eta}P_{\bar\eta}:H^*(C_{\bar\eta},\mathbf Q_\ell(i))\longrightarrow H^{*+2(m-1)}(X_{\bar\eta},\mathbf Q_\ell(i+m-1)).
\end{equation}
考虑\Term{algebraic-correspondence}{代数对应}
\begin{equation}\tag{4.4.10}
\begin{aligned}
R_{\bar\eta}&=P_{\bar\eta}-m^*Q_{\bar\eta}=(\operatorname{id}-m^*V_{\bar\eta})P_{\bar\eta}:\\
&H^*(C_{\bar\eta},\mathbf Q_\ell(i))\longrightarrow H^{*+2(m-1)}(X_{\bar\eta},\mathbf Q_\ell(i+m-1)).
\end{aligned}
\end{equation}
由于 $V_{\bar\eta}$ 是 $m^*$ 的\Term{left-inverse}{左逆}，$\operatorname{id}-m^*V_{\bar\eta}$ 是 $H^*(X_{\bar\eta},-)$ 中的一个\Term{projector}{投影算子}，其像是
\[\operatorname{Im}\bigl(m^*:H^*(X\otimes_k k(\bar\eta),-)\longrightarrow H^*(X_{\bar\eta},-)\bigr)\]
的一个\Term{direct-complement}{直和补空间}。由假设（4.1.2），\Term{homomorphism}{同态}
\begin{equation}\tag{4.4.11}
R_{\bar\eta}:H^1(C_{\bar\eta},\mathbf Q_\ell(i))\longrightarrow H^{2m-1}(X_{\bar\eta},\mathbf Q_\ell(i+m-1))
\end{equation}
的像包含于 $\operatorname{Im}m^*$，故得到：

\noindent（4.4.12）\Term{homomorphism}{同态}（4.4.11）为零。

现在选取一个\Term{finite-extension}{有限扩张}
\[L/k(\eta),\]
使\Term{object}{对象}
\[C_{\bar\eta},\ t_1(\bar\eta),\ t_2(\bar\eta),\ Q_{\bar\eta},\ R_{\bar\eta}\]
都已经在它上面有定义；并取 $\mathbf P^1$ 在 $L$ 中的\Term{normalization}{正规化}的一个\Term{nonempty}{非空}\Term{open-subscheme}{开子概形} $U$。把 $U$ 取到足够小，就存在一条\Term{proper}{固有}\Term{smooth}{光滑}、\Term{fibre}{纤维}\Term{geometrically-connected}{几何连通}的\Term{relative-curve}{相对曲线}
\SourcePage{369}
\[C\longrightarrow U,\]
它的\Term{geometric-generic-fibre}{几何泛纤维}为 $C_{\bar\eta}$；并有两个\Term{section}{截面}
\[t_i:U\longrightarrow C,\qquad i=1,2,\]
在 $\bar\eta$ 处的值是 $t_i(\bar\eta)$。设
\[\rho:\widetilde X\longrightarrow\mathbf P^1\]
是\Term{fibre}{纤维}为\Term{lefschetz-pencil}{铅笔}各\Term{section}{截面}的\Term{projection}{投影}（第 XVIII 讲 3.1.4），并令
\[\widetilde X_U=\widetilde X\times_{\mathbf P^1}U,\]
且令
\[M:\widetilde X_U\hookrightarrow X\times U\]
为 $\widetilde X_U$ 到 $X\times U$ 的\Term{canonical}{典范}\Term{immersion}{浸入}（把后者看作 $U$ 上取值为 $X$ 的\Term{constant-family}{常值族}），以及
\[z_U=M^*\operatorname{pr}_1^*(z)\]
为 $z$ 在 $\widetilde X_U$ 上诱导的\Term{algebraic-cycle}{代数圈}。把 $U$ 取到足够小，就存在 $C\times_U\widetilde X_U$ 上的\Term{algebraic-cycle}{代数圈} $R$，其在 $C_{\bar\eta}\times X_{\bar\eta}$ 上的值为 $R_{\bar\eta}$；还存在 $C\times X\times U$ 上的\Term{algebraic-cycle}{代数圈} $Q$，其在 $C_{\bar\eta}\times(X\otimes_k k(\bar\eta))$ 上的值为 $Q_{\bar\eta}$，并且满足
\begin{equation}\tag{4.4.13}
z_U=R(t_1)-R(t_2)+(\operatorname{id}\times M)^*[Q(t_1)-Q(t_2)].
\end{equation}
由 \Term{griffiths-homomorphism}{Griffiths 同态}“模去\Term{fixed-part}{固定部分}”（第 XVIII 讲 5.8.7）的\Term{injectivity}{单射性}，
\SourcePage{370}
\[
\operatorname{Prim}^{2m}(X,\mathbf Q_\ell(n))\xrightarrow{\operatorname{grif}}
\frac{H^1(U,R^{2m-1}\rho_{U*}\mathbf Q_\ell(m))}{M^*H^1(U,H^{2m-1}(X,\mathbf Q_\ell(m)))},
\]
要证明\Term{theorem}{定理}，只需证明 $z_U$ 在
\[
u_\rho:\operatorname{Prim}^{2m}(\widetilde X_U/U,\mathbf Q_\ell(m))\longrightarrow H^1(U,R^{2m-1}\rho_{U*}\mathbf Q_\ell(m))
\]
下的像落在 $M^*\bigl(H^1(U,H^{2m-1}(X,\mathbf Q_\ell(m)))\bigr)$ 中。

由\Term{resolution}{分解}式（4.4.13），以及（1.4.2）的\Term{commutativity}{交换性}，即
\[
\begin{tikzcd}[column sep=small]
\operatorname{Prim}^2(C/U,\mathbf Q_\ell(1)) \arrow[r,"Q"] & \operatorname{Prim}^{2m}(X\times U/U,\mathbf Q_\ell(n)) \arrow[r,"M^*"] \arrow[d,"u_{\operatorname{pr}_2}"'] & \operatorname{Prim}^{2m}(\widetilde X_U/U,\mathbf Q_\ell(m)) \arrow[d,"u_\rho"]\\
& H^1(U,H^{2m-1}(X,\mathbf Q_\ell(m))) \arrow[r,"M^*"'] & H^1(U,R^{2m-1}\rho_{U*}\mathbf Q_\ell(m)),
\end{tikzcd}
\]
只需证明
\[
u_\rho(R(t_1)-R(t_2))=0\quad\text{于 }H^1(U,R^{2n-1}\rho_{U*}\mathbf Q_\ell(n)),
\]
而由 1.7，这正是（4.3.12）的推论。

\SGAsection{5}{Griffiths 群}
\textbf{5.0。} 设 $S$ 在\Term{algebraically-closed}{代数闭}域 $k$ 上\Term{proper}{固有}且\Term{smooth}{光滑}。用
\SourcePage{371}

\noindent（5.0.1）$Z^i(S)$ 表示 $S$ 上\Term{codimension}{余维}为 $i$ 的\Term{algebraic-cycle}{代数圈}群；

\noindent（5.0.2）$Z^i_{\mathrm{alg}}(S)$ 表示 $Z^i(S)$ 中由\Term{algebraic-equivalence}{代数等价}于零的圈组成的\Term{subgroup}{子群}；

并对每个素数 $\ell\ne\operatorname{char}(k)$，用

\noindent（5.0.3）$Z^i_{\ell\text{-}\mathrm{coh}}(S)$ 表示 $Z^i(S)$ 中由其在 $H^{2i}(S,\mathbf Q_\ell(i))$ 中的\Term{cohomology-class}{上同调类}为零的圈组成的\Term{subgroup}{子群}。

（人们猜想 [4] 这个群与 $\ell$ 无关；至少在 $\operatorname{char}k=0$ 时这是成立的。）置
\begin{equation}\tag{5.0.4}
\operatorname{Grif}^i_\ell(S)=Z^i_{\ell\text{-}\mathrm{coh}}(S)/Z^i_{\mathrm{alg}}(S).
\end{equation}
若 $S$ 还是\Term{projective}{射影}的，并且配备一个\Term{projective-immersion}{射影浸入}，则置

\noindent（5.0.5）$Z^{i,\mathrm{prim}}(S)$ 为 $Z^i(S)$ 中由\Term{primitive-cycle}{本原圈}（即其\Term{cohomology-class}{上同调类}为\Term{primitive-class}{本原的}）组成的\Term{subgroup}{子群}；

\noindent（5.0.6）$\operatorname{Prim}^{2i}_{\mathrm{alg}}(S,\mathbf Q_\ell(i))$ 为 $H^{2i}(S,\mathbf Q_\ell(i))$ 中由\Term{primitive-cycle}{本原代数圈}的\Term{cohomology-class}{上同调类}生成的 $\mathbf Q$-\Term{vector-subspace}{向量子空间}。

\begin{SGAresult}{推论}{5.1}
在 4.1 的假设下，有
\begin{equation}\tag{5.1.1}
\dim_{\mathbf Q}\bigl(\operatorname{Grif}^{n}(X_{\bar\eta})\otimes_{\mathbf Z}\mathbf Q\bigr)\geq\dim_{\mathbf Q}\bigl(\operatorname{Prim}^{2n}_{\mathrm{alg}}(X,\mathbf Q_\ell(n))\bigr).
\end{equation}
\end{SGAresult}
\begin{proof}
由 4.1.5，复合映射
\[
\mathbf Q\otimes_{\mathbf Z}Z^{n,\mathrm{prim}}(X)\xrightarrow{\text{限制到 }X_{\bar\eta}}\mathbf Q\otimes_{\mathbf Z}Z^n_{\ell\text{-}\mathrm{coh}}(X_{\bar\eta})\longrightarrow\operatorname{Grif}^n(X_{\bar\eta})
\]
的核
\SourcePage{372}
包含于 $Z^n_{\ell\text{-}\mathrm{coh}}(X)\otimes_{\mathbf Z}\mathbf Q$，因此 $\operatorname{Prim}^{2n}_{\mathrm{alg}}(X,\mathbf Q_\ell(n))$ 是 $\operatorname{Grif}^n(X_{\bar\eta})\otimes_{\mathbf Z}\mathbf Q$ 的一个\Term{subgroup}{子群}的商。这就完成了证明。
\end{proof}
\textbf{注记 5.2。} 若 $X$ 是 $\mathbf P^{2m+1}$ 中的一个\Term{quadric-hypersurface}{二次超曲面}，则
\[\dim_{\mathbf Q_\ell}\bigl(\operatorname{Prim}^{2m}_{\mathrm{alg}}(X,\mathbf Q_\ell(n))\bigr)=1.\]
特别地（4.2），若 $X$ 是 $\mathbf P^5$ 中的一个\Term{quadric-hypersurface}{二次超曲面}，则用\Term{degree}{次数} $>5$ 的\Term{generic-hypersurface}{泛超曲面}截 $X$ 所得的\Term{section}{截面}包含一条\Term{curve}{曲线}，它的 $\mathbf Q_\ell$-\Term{cohomology-class}{上同调类}为零，但其任何倍数都不\Term{algebraic-equivalence}{代数等价}于零。

\SGAsection{}{参考文献}
\begin{enumerate}
\item[{[1]}] Griffiths, P.，\emph{某些有理积分的周期 III}（On the Periods of Certain Rational Integrals III），即将刊于 \emph{Annals of Mathematics}。
\item[{[2]}] Grothendieck, A.，\emph{Brauer 群 III}（Le Groupe de Brauer III，特别参见第 9、10 节），载 \emph{\Term{scheme}{概形}\Term{cohomology}{上同调}十讲}（Dix Exposés sur la Cohomologie des Schémas），North Holland，1968。
\item[{[3]}] Jouanolou, J. P.，\emph{$\ell$-进\Term{cohomology}{上同调}}（Cohomologie $\ell$-adique），学位论文，巴黎，1969。
\item[{[4]}] Kleiman, S.，\emph{\Term{algebraic-cycle}{代数圈}与 Weil 猜想}（Algebraic Cycles and the Weil Conjectures），载 \emph{\Term{scheme}{概形}\Term{cohomology}{上同调}十讲}（Dix Exposés sur la Cohomologie des Schémas），North Holland，1968。
\item[{[5]}] Manin, J.，\emph{\Term{function-field}{函数域}上\Term{algebraic-curve}{代数曲线}的\Term{rational-point}{有理点}}（Rational Points of Algebraic Curves over Function Fields），\emph{Translations Amer. Math. Soc.}\SGAsourcegap{原书文献 [5] 到期刊简称即止，未提供卷次、出版年或页码。}
\end{enumerate}
